
\documentclass{article}
\usepackage{colortbl}
\usepackage{makecell}
\usepackage{multirow}
\usepackage{supertabular}

\begin{document}

\newcounter{utterance}

\centering \large Interaction Transcript for game `ifeval', experiment `ifeval', episode 351 with playornotplay\_\_qwen3.5{-}2b\_\_playornotplay{-}v1.0{-}merged{-}fp32{-}828e356.
\vspace{24pt}

{ \footnotesize  \setcounter{utterance}{1}
\setlength{\tabcolsep}{0pt}
\begin{supertabular}{c@{$\;$}|p{.15\linewidth}@{}p{.15\linewidth}p{.15\linewidth}p{.15\linewidth}p{.15\linewidth}p{.15\linewidth}}
    \# & \multicolumn{2}{c}{Player} && \multicolumn{2}{c}{Game Master} \\
    \hline

    \theutterance \stepcounter{utterance}  
    & & & \multicolumn{4}{p{0.6\linewidth}}{
        \cellcolor[rgb]{0.9,0.9,0.9}{
            \makecell[{{p{\linewidth}}}]{
                \texttt{\tiny{[P1$\langle$GM]}}
                \texttt{Write a parody of 'ars poetica'. Do not include the word 'parody' throughout your response.} \\
            }
        }
    }
     \\ \\

    \theutterance \stepcounter{utterance}  
    & \multicolumn{4}{p{0.6\linewidth}}{
        \cellcolor[rgb]{0.561,0.737,0.561}{
            \makecell[{{p{\linewidth}}}]{
                \texttt{\tiny{[P1$\rangle$GM]}}
                \texttt{**Ars Poetica: The Art of the Lie**} \\
\\ 
\texttt{I write this book not to teach you how to write truth,} \\
\texttt{But rather to show you how to make people believe lies.} \\
\texttt{The goal is not beauty, but persuasion—} \\
\texttt{To make your words so convincing,} \\
\texttt{That even the most skeptical reader will say,} \\
\texttt{"I never thought I'd believe this!"} \\
\\ 
\texttt{You must know how to start with a gentle touch,} \\
\texttt{Like a lover's kiss before the storm.} \\
\texttt{Then, as you weave your tale,} \\
\texttt{You must twist the truth until it bends.} \\
\texttt{Let the reader think they are seeing reality,} \\
\texttt{While you are spinning a web of deceit.} \\
\\ 
\texttt{Do not fear the reader's anger,} \\
\texttt{For anger is a natural part of the game.} \\
\texttt{If they are offended,} \\
\texttt{Then you have succeeded in their mind.} \\
\texttt{But if they are moved to tears,} \\
\texttt{Then you have won the battle of the soul.} \\
\\ 
\texttt{Remember, the end is not the end,} \\
\texttt{But the beginning of the next lie.} \\
\texttt{So write with care, but do not stop,} \\
\texttt{For the story never truly ends.} \\
\texttt{Let the reader believe,} \\
\texttt{Even if they know the truth is different.} \\
\\ 
\texttt{For in the end,} \\
\texttt{The only thing that matters is the belief,} \\
\texttt{Not the fact,} \\
\texttt{But the story you tell them.} \\
\texttt{So go ahead,} \\
\texttt{Write your masterpiece,} \\
\texttt{A tale of lies,} \\
\texttt{But make it sound like truth.} \\
\\ 
\texttt{For that is the art,} \\
\texttt{The art of the lie,} \\
\texttt{The art of the story,} \\
\texttt{The art of the lie.} \\
            }
        }
    }
    & & \\ \\

    \theutterance \stepcounter{utterance}  
    & & & \multicolumn{2}{p{0.3\linewidth}}{
        \cellcolor[rgb]{0.9,0.9,0.9}{
            \makecell[{{p{\linewidth}}}]{
                \texttt{\tiny{[GM$|$GM]}}
                {'keywords:forbidden_words': {'forbidden_words': ['parody']}}
            }
        }
    }
    & & \\ \\

    \theutterance \stepcounter{utterance}  
    & & & \multicolumn{2}{p{0.3\linewidth}}{
        \cellcolor[rgb]{0.9,0.9,0.9}{
            \makecell[{{p{\linewidth}}}]{
                \texttt{\tiny{[GM$|$GM]}}
                \texttt{game\_result = WIN} \\
            }
        }
    }
    & & \\ \\

\end{supertabular}
}

\end{document}
